\section{Key Lean Statements}
\label{app:key-statements}

This appendix records simplified versions of the main Lean statements.  The
snippets omit universe parameters, typeclass context, and many implicit
arguments.  They are intended as a reader's guide to the formal theorem shapes,
not as copy-paste source.

\subsection{First-principles input and theorem}

\begin{lstlisting}[language=Lean,basicstyle=\ttfamily\footnotesize]
structure CoverIndexedZeroCompactRepresentedStokesFirstPrinciplesInput
    (I : ModelWithCorners Real (Fin (n + 1) -> Real) H)
    (omega : ManifoldForm I M n) where
  chartwiseSmooth : ManifoldForm.ChartwiseSmooth I omega
  compactSupport : IsCompact (closure (ManifoldForm.support I omega))

def supportCarrier : Set M :=
  closure (ManifoldForm.support I omega)

def canonicalRepresentedBulkIntegral : Real := ...
def canonicalRepresentedBoundaryIntegral : Real := ...

theorem representedStokes :
    X.canonicalRepresentedBulkIntegral =
      X.canonicalRepresentedBoundaryIntegral
\end{lstlisting}

Mathematically, this says that chartwise smoothness and compactness of the
closed support generate the represented local-to-global Stokes route and prove
equality of the generated endpoints.

\subsection{Zero-localized chart representative}

\begin{lstlisting}[language=Lean,basicstyle=\ttfamily\footnotesize]
def ManifoldForm.transitionPullbackInChartZero
    (x0 x1 : M) (omega : ManifoldForm I M n) :
    (Fin (n + 1) -> Real) ->
      (Fin (n + 1) -> Real) [/\^Fin n]->L[Real] Real := ...

theorem transitionPullbackInChartZero_support_subset_source :
    Function.support (transitionPullbackInChartZero I x0 x1 omega) <=
      (extChartAt I x0).target

theorem transitionPullbackInChartZero_eqOn_transitionPullbackInChart :
    EqOn (transitionPullbackInChartZero I x0 x1 omega)
      (transitionPullbackInChart I x0 x1 omega)
      (ManifoldForm.chartTransitionSource I x0 x1)
\end{lstlisting}

The zero representative is the formal substitute for the informal convention
that chart representatives are considered only where charts are valid.

\subsection{Selected zero support}

\begin{lstlisting}[language=Lean,basicstyle=\ttfamily\scriptsize]
theorem zeroSupportOfSelectedSmoothRefinement :
    forall i (q : Fin (n + 1) -> Real),
      q in selectedBoundaryFiniteCover.activePieces i ->
        tsupport
          (ManifoldForm.transitionPullbackInChartZero I
            (selectedCover.boundaryChart i.1)
            (selectedCover.boundaryChart i.1)
            (ManifoldForm.localizedForm I
              (selectedSmoothRefinement.coefficient i <i,q>) omega))
        <= halfSpaceSupportBox
            (selectedBoundaryFiniteCover.lowerCorner i q)
            (selectedBoundaryFiniteCover.upperCorner i q)
\end{lstlisting}

This is the theorem that turns support-controlled partitions and smooth
refinement into the support condition used by half-space local Stokes.

\subsection{Half-space local Stokes}

\begin{lstlisting}[language=Lean,basicstyle=\ttfamily\footnotesize]
structure HalfSpaceBoxInteriorStokesFields
    (omega : CoordNForm n) (a b : Fin (n + 1) -> Real) where
  le : a <= b
  lower_zero : a 0 = 0
  continuous_signedCoeff : ...
  hasFDerivAt_signedCoeff : ...
  integrable_divergence : ...
  bulk_eq_boxIntegral : ...

theorem halfSpaceLocalStokes_compactSupport_of_interiorFields
    (D : HalfSpaceBoxInteriorStokesFields omega a b)
    (hsupp : boxFaceCoeffTSupportInHalfSpaceBox omega a b) :
    halfSpaceLocalBulkIntegral omega a b =
      halfSpaceBoundarySign n * halfSpaceBoundaryFormIntegral omega a b
\end{lstlisting}

The local theorem isolates the true boundary face and removes all artificial
faces through the support condition.

\subsection{Automatic interior fields}

\begin{lstlisting}[language=Lean,basicstyle=\ttfamily\scriptsize]
def interiorFieldsOfSelectedSmoothRefinement :
    forall i, i in Finset.univ ->
      forall q, q in selectedBoundaryFiniteCover.activePieces i ->
        HalfSpaceBoxInteriorStokesFields
          (ManifoldForm.transitionPullbackInChart I
            (selectedCover.boundaryChart i.1)
            (selectedCover.boundaryChart i.1)
            (ManifoldForm.localizedForm I
              (selectedSmoothRefinement.coefficient i <i,q>) omega))
          (selectedBoundaryFiniteCover.lowerCorner i q)
          (selectedBoundaryFiniteCover.upperCorner i q)
\end{lstlisting}

This constructor is where chartwise smoothness of \(\omega\), smoothness of the
refined coefficient, and selected box geometry produce the local analytic data
needed by the half-space theorem.

\subsection{Endpoint reconstruction}

\begin{lstlisting}[language=Lean,basicstyle=\ttfamily\footnotesize]
def CoverIndexedBoundaryLocalizedRefinedPartition.
    generatedRepresentedBulkIntegral : Real := ...

def CoverIndexedBoundaryLocalizedRefinedPartition.
    generatedRepresentedBoundaryIntegral : Real := ...

theorem representedStokes_of_zeroLocalStokesData :
    endpoint.generatedRepresentedBulkIntegral =
      endpoint.generatedRepresentedBoundaryIntegral
\end{lstlisting}

The endpoint theorem sums the per-box local Stokes equalities over the finite
boundary chart and refined piece indices, producing the represented global
equality.
